% Generated draft: requires mathematical and source review.
\documentclass[11pt]{article}
\usepackage[margin=25mm]{geometry}
\usepackage{fontspec}
\setmainfont{Latin Modern Roman}
\usepackage{amsmath,amssymb,mathtools,unicode-math}
\setmathfont{Latin Modern Math}
\usepackage{enumitem,needspace,xurl,hyperref}
\setlength{\parindent}{0pt}
\setlength{\parskip}{5pt}
\setlength{\emergencystretch}{3em}
\newcommand{\N}{\mathbb{N}}
\newcommand{\Z}{\mathbb{Z}}
\newcommand{\Q}{\mathbb{Q}}
\newcommand{\R}{\mathbb{R}}
\newcommand{\C}{\mathbb{C}}
\newcommand{\F}{\mathbb{F}}
\newcommand{\A}{\mathbb{A}}
\newcommand{\PP}{\mathbb{P}}
\newcommand{\E}{\mathbb{E}}
\newcommand{\eps}{\varepsilon}
\newcommand{\dd}{\,\mathrm{d}}
\begin{document}
\section*{UnsolvedMath ID: 30004753}
\textbf{Source problem code:} catalogue entry
\subsection{Definitions and mathematical statement}
Let \(p\) be a prime. Regard \(\F_{p^2}\) as a two-dimensional vector space over \(\F_p\). After choosing a basis, multiplication by its nonzero elements identifies \(\F_{p^2}^{\times}\) with a nonsplit Cartan subgroup \(C_{\mathrm{ns}}(p)\subseteq\operatorname{GL}_2(\F_p)\). Define its normalizer by
\[
N_{\mathrm{ns}}(p)=\{g\in\operatorname{GL}_2(\F_p):gC_{\mathrm{ns}}(p)g^{-1}=C_{\mathrm{ns}}(p)\},
\]
and put
\[
\Gamma_{\mathrm{ns}}^+(p)=\{\gamma\in\operatorname{SL}_2(\Z):\gamma\bmod p\in N_{\mathrm{ns}}(p)\}.
\]
Use the standard interpretation of \(X(\Gamma_{\mathrm{ns}}^+(p))\) as the smooth projective modular curve over \(\Q\) associated with \(N_{\mathrm{ns}}(p)\), often denoted \(X_{\mathrm{ns}}^+(p)\). Its complex points are described by the compactified quotient
\[
X_{\mathrm{ns}}^+(p)(\C)\simeq\Gamma_{\mathrm{ns}}^+(p)\backslash\bigl(\mathfrak H\cup\PP^1(\Q)\bigr),
\qquad \mathfrak H=\{z\in\C:\operatorname{Im}(z)>0\}.
\]
The rational structure here is the standard one attached to the normalizer; the complex quotient alone does not specify a model over \(\Q\).

\textbf{Problem.} Determine the complete set
\[
X_{\mathrm{ns}}^+(p)(\Q)
\]
as \(p\) ranges over the prime levels intended by the source. Under the projective \(X\)-convention above, this includes any rational cusps.

\textbf{Unresolved scope qualifications.} The supplied background describes a question about prime levels, but neither statement specifies a lower bound, exceptional primes, or a more restricted range of \(p\). Thus the unrestricted prime family is the natural reading, while any intended restriction requires source confirmation. The compactification and rational model above make explicit the standard interpretation of the notation; the supplied statement does not separately describe them or exclude cusps.
\subsection{Short English statement}
Determine all rational points on the modular curves attached to normalizers of nonsplit Cartan subgroups at prime level \(p\). The supplied record does not specify a restricted range of primes. With the standard projective interpretation of \(X\), rational cusps are included. The record labels the problem open; that label is supplied curation, not an independently verified conclusion.
\subsection{Sources}
\begin{itemize}
\item Local UnsolvedMath registry: permanent ID \(30004753\), title \emph{Rational Points on Nonsplit Cartan Modular Curves}, exact source problem code \texttt{catalogue entry}; collection: \emph{Oberwolfach Reports Open Problems}. The supplied record gives proposed year 2021 and does not name a proposer.
\item Supplied report citation: \emph{Explicit Methods in Number Theory (2022)}, Oberwolfach Reports, DOI 10.4171/owr/2021/34, pp.\ 1805--1858. \url{https://doi.org/10.4171/owr/2021/34}. The background describes Oberwolfach Reports as documenting mathematical workshops and their open-problem discussions.
\item Supplied literature source, labelled Rouse--Sutherland--Zureick-Brown: \url{https://arxiv.org/abs/2106.11141}.
\item Supplied literature source, labelled Bilu--Parent: \url{https://annals.math.princeton.edu/2011/173-1/p13}.
\item Literature assessment dated exactly \texttt{2026-08-22}: the supplied curation labels the problem open and reports that rational points on these plus nonsplit Cartan modular curves are known in substantial special cases, while a uniform determination for the prime levels in question remains open. This is supplied dated curation derived from the OWR material, not an independently verified literature assessment; its time-sensitive status and cited claims require independent verification.
\end{itemize}
\end{document}
